\section{How to add and evaluate an observer}
\label{sec:artifact}

ObserverBench uses the same four-step workflow for every submission:

\begin{enumerate}[leftmargin=1.4em,itemsep=0.25em,topsep=0.3em]
  \item \textbf{Choose a fixed task.} Its TaskCard states what is predicted,
  what the observer may read, which actions are allowed, how actions are chosen,
  which data are held out, and how success is scored.
  \item \textbf{Describe and submit the observer.} Its ObserverCard records the
  features it uses, its access, fitting procedure, implementation, and known
  failures. The method may be submitted as code or as a table of predictions.
  \item \textbf{Run the common evaluator.} The task applies the same held-out
  split, action rule, and budget to every observer. The observer cannot replace
  the controller or reveal the test targets.
  \item \textbf{Report both levels of performance.} The result separates
  prediction or ranking quality from the loss caused by the chosen action. It
  also records what information and extra model calls the observer required.
\end{enumerate}

The public release will turn each completed evaluation into a task-specific
scorecard and compare it with the bundled baseline results. A leaderboard
position is meaningful only among results that use the same task version,
information boundary, action rule, and budget. ObserverBench therefore does
not combine unlike settings into one overall rank.

The release supports three task types:

\begin{itemize}
  \item \textbf{Control} separates the state estimator, edit direction,
  controller, and task. A submission can replace the estimator or direction
  while the model, starting states, controller settings, and metrics remain
  fixed. The analytic control task provides the tested external adapter.
  \item \textbf{Finite-effect prediction} provides fixed measurements and asks
  for the effects of held-out interventions. The task owns the measurement
  budget, split, targets, action rule, and metrics. A table submission does not
  need the original model.
  \item \textbf{Safety triage} asks for one risk score per held-out request or
  trajectory. The evaluator converts scores into allow, block, escalate, or
  audit actions under fixed budgets and reports realized violations, tail loss,
  false interventions, and clean utility.
\end{itemize}

The versioned interfaces are \texttt{observerbench.api.v0},
\texttt{observerbench.effect\_prediction.v0}, and
\texttt{observerbench.safety\_protocol.v0}. Public safety-task packs describe
the threat model, actions, loss, budgets, training measurements, and held-out
queries. They omit the held-out answers, which the evaluator stores separately
and links to the public task by hash. This lets researchers add data-only safety
tasks without changing the workbench.

The bundled GPT-2 IOI and Qwen induction-copy registries show the finite-effect
workflow at several measurement budgets. Both expose frozen tables, check the
declared model and split, and use separate prompts for measurements and held-out
targets. Outside predictors therefore need neither model inference nor the
original experiment code; per-head or additive ridge observers provide the
baselines.

\subsection{A four-step check for safety monitors}

Before interpreting an internal safety signal, ObserverBench asks:

\begin{enumerate}[leftmargin=1.4em,itemsep=0.2em,topsep=0.25em]
  \item \textbf{Does the model output already suffice?} Read the relevant output
  distribution rather than one sampled answer. If it is adequate, stop.
  \item \textbf{Does the input already suffice?} Train a text or input-feature
  baseline with the same supervision. If it matches the internal probe, an
  activation-specific explanation is unnecessary.
  \item \textbf{Does the signal remain without a request for a verdict?} Measure
  the same examples in a neutral context to separate input processing from a
  signal created or reshaped by the monitoring prompt.
  \item \textbf{Does it work under the intended policy?} Test frozen scores at
  the attack rate, intervention budget, deployment groups, and loss the
  controller will face.
\end{enumerate}

These checks locate where useful information becomes available. They do not by
themselves establish deception, intent, or hidden knowledge. Their purpose is a
fair comparison: every observer receives the same task and faces the same action
rule.
